% Shared preamble of the proof overview documents  <Folder>/<Name>_overview.tex.
% Each overview starts with
%     % Shared preamble of the proof overview documents  <Folder>/<Name>_overview.tex.
% Each overview starts with
%     % Shared preamble of the proof overview documents  <Folder>/<Name>_overview.tex.
% Each overview starts with
%     \input{../../scripts/overview_preamble.tex}
% (for ThinMatching/: \input{../scripts/overview_preamble.tex}) and is compiled from its own folder by
% scripts/build_overviews.py.
\documentclass[10pt,a4paper]{article}
\usepackage[T1]{fontenc}
\usepackage[utf8]{inputenc}
\usepackage[margin=2.2cm]{geometry}
\usepackage{amsmath,amssymb,amsthm}
\usepackage{array,booktabs,longtable}
\usepackage[dvipsnames]{xcolor}
\usepackage{enumitem}
\usepackage{tikz}
\usetikzlibrary{arrows.meta,positioning,shapes.multipart,fit,calc}
\usepackage[hidelinks]{hyperref}
\setlist{itemsep=2pt,topsep=4pt}
\renewcommand{\arraystretch}{1.2}

% Lean names, files and paper references
\newcommand{\lean}[1]{\textcolor{blue!50!black}{\texttt{\detokenize{#1}}}}
\newcommand{\file}[1]{\textcolor{OliveGreen!70!black}{\texttt{\detokenize{#1}}}}
\newcommand{\paperref}[1]{\textsf{#1}}
\newcommand{\hyp}[1]{\textcolor{BrickRed}{\textbf{#1}}}

% Boxes for the dependency diagrams
\tikzset{
  res/.style={draw, rounded corners, fill=blue!6, align=center, inner sep=4pt, font=\small},
  main/.style={res, fill=blue!18, very thick},
  cited/.style={res, fill=red!10, dashed},
  shared/.style={res, fill=gray!12},
  dep/.style={-{Stealth[length=2.2mm]}, thick},
}

\newcommand{\overviewtitle}[2]{%
  \begin{center}{\Large\bfseries #1}\\[4pt]{\small Proof overview of the Lean~4/Mathlib formalisation}\\[2pt]
  {\small #2}\end{center}}

% Standard paragraph describing the paper
\newcommand{\paperversion}{\noindent\textbf{Paper.} J.~Breitung, T.~Gervens, M.~Grohe.
  \emph{Graph Similarity Meets Optimal Transport}.\\ File \file{Graph_Similarity_Meets_Optimal_Transport.pdf} in the
  repository; statement, equation, page and line numbers refer to this PDF. \file{docs/NUMBERING.md} lists the Lean
  files and declarations of all numbered statements.\par\smallskip}

% (for ThinMatching/: \input{../scripts/overview_preamble.tex}) and is compiled from its own folder by
% scripts/build_overviews.py.
\documentclass[10pt,a4paper]{article}
\usepackage[T1]{fontenc}
\usepackage[utf8]{inputenc}
\usepackage[margin=2.2cm]{geometry}
\usepackage{amsmath,amssymb,amsthm}
\usepackage{array,booktabs,longtable}
\usepackage[dvipsnames]{xcolor}
\usepackage{enumitem}
\usepackage{tikz}
\usetikzlibrary{arrows.meta,positioning,shapes.multipart,fit,calc}
\usepackage[hidelinks]{hyperref}
\setlist{itemsep=2pt,topsep=4pt}
\renewcommand{\arraystretch}{1.2}

% Lean names, files and paper references
\newcommand{\lean}[1]{\textcolor{blue!50!black}{\texttt{\detokenize{#1}}}}
\newcommand{\file}[1]{\textcolor{OliveGreen!70!black}{\texttt{\detokenize{#1}}}}
\newcommand{\paperref}[1]{\textsf{#1}}
\newcommand{\hyp}[1]{\textcolor{BrickRed}{\textbf{#1}}}

% Boxes for the dependency diagrams
\tikzset{
  res/.style={draw, rounded corners, fill=blue!6, align=center, inner sep=4pt, font=\small},
  main/.style={res, fill=blue!18, very thick},
  cited/.style={res, fill=red!10, dashed},
  shared/.style={res, fill=gray!12},
  dep/.style={-{Stealth[length=2.2mm]}, thick},
}

\newcommand{\overviewtitle}[2]{%
  \begin{center}{\Large\bfseries #1}\\[4pt]{\small Proof overview of the Lean~4/Mathlib formalisation}\\[2pt]
  {\small #2}\end{center}}

% Standard paragraph describing the paper
\newcommand{\paperversion}{\noindent\textbf{Paper.} J.~Breitung, T.~Gervens, M.~Grohe.
  \emph{Graph Similarity Meets Optimal Transport}.\\ File \file{Graph_Similarity_Meets_Optimal_Transport.pdf} in the
  repository; statement, equation, page and line numbers refer to this PDF. \file{docs/NUMBERING.md} lists the Lean
  files and declarations of all numbered statements.\par\smallskip}

% (for ThinMatching/: \input{../scripts/overview_preamble.tex}) and is compiled from its own folder by
% scripts/build_overviews.py.
\documentclass[10pt,a4paper]{article}
\usepackage[T1]{fontenc}
\usepackage[utf8]{inputenc}
\usepackage[margin=2.2cm]{geometry}
\usepackage{amsmath,amssymb,amsthm}
\usepackage{array,booktabs,longtable}
\usepackage[dvipsnames]{xcolor}
\usepackage{enumitem}
\usepackage{tikz}
\usetikzlibrary{arrows.meta,positioning,shapes.multipart,fit,calc}
\usepackage[hidelinks]{hyperref}
\setlist{itemsep=2pt,topsep=4pt}
\renewcommand{\arraystretch}{1.2}

% Lean names, files and paper references
\newcommand{\lean}[1]{\textcolor{blue!50!black}{\texttt{\detokenize{#1}}}}
\newcommand{\file}[1]{\textcolor{OliveGreen!70!black}{\texttt{\detokenize{#1}}}}
\newcommand{\paperref}[1]{\textsf{#1}}
\newcommand{\hyp}[1]{\textcolor{BrickRed}{\textbf{#1}}}

% Boxes for the dependency diagrams
\tikzset{
  res/.style={draw, rounded corners, fill=blue!6, align=center, inner sep=4pt, font=\small},
  main/.style={res, fill=blue!18, very thick},
  cited/.style={res, fill=red!10, dashed},
  shared/.style={res, fill=gray!12},
  dep/.style={-{Stealth[length=2.2mm]}, thick},
}

\newcommand{\overviewtitle}[2]{%
  \begin{center}{\Large\bfseries #1}\\[4pt]{\small Proof overview of the Lean~4/Mathlib formalisation}\\[2pt]
  {\small #2}\end{center}}

% Standard paragraph describing the paper
\newcommand{\paperversion}{\noindent\textbf{Paper.} J.~Breitung, T.~Gervens, M.~Grohe.
  \emph{Graph Similarity Meets Optimal Transport}.\\ File \file{Graph_Similarity_Meets_Optimal_Transport.pdf} in the
  repository; statement, equation, page and line numbers refer to this PDF. \file{docs/NUMBERING.md} lists the Lean
  files and declarations of all numbered statements.\par\smallskip}

\begin{document}
\overviewtitle{Proposition 5: the lower bound $4/3$ for the edit distance}{Folder \file{GsimOT/Proposition6_EditLowerBound}}
\paperversion

\section{Statement}

\paragraph{Paper.} Proposition~5 (Section~3.1, p.~6, l.~221--223; proof in App.~C.1, pp.~21--24,
l.~675--758): \emph{For every odd $N\ge3$, there are graphs $G$ and $H$ with $d_{\mathrm{ed}}(G,H)\ge4N-2$
and $d^{\mathrm{OT}}_{\mathrm{ed}}(G,H)\le3N+1$. Thus, for every $\varepsilon>0$, there is a sufficiently
large $N$, such that $\delta_{\mathrm{ed}}(G,H)\ge(\frac43-\varepsilon)\,\delta^{\mathrm{OT}}_{\mathrm{ed}}(G,H)$.}
The second sentence gives Theorem~3~(1) (p.~4, l.~124).

\paragraph{Lean.} The graphs are $G=$ \lean{GN N} and $H=$ \lean{HN N} on the common vertex set
\lean{VN N} (\file{Construction.lean}). All statements are in the namespace \lean{GsimOT.Prop6}.
\begin{itemize}
\item \lean{proposition6_dEdit}: if $N$ is odd and $N\ge3$, then
  $4N-2\le d_{\mathrm{ed}}(G,H)$ and $d^{\mathrm{OT}}_{\mathrm{ed}}(G,H)\le3N+1$ \quad (the paper's first sentence).
\item \lean{proposition6}: the same hypotheses give $4N-2\le C_{\mathrm{ed}}(\pi)$ for every bijection $\pi$ and
  an explicit $Q\in\mathrm{DS}(V,W)$ with $C_{\mathrm{ed}}(Q)\le3N+1$ (the form with a witness).
\item \lean{proposition6_delta_of_le}: if $N$ is odd, $N\ge3$ and $\frac{10}{3(3N+1)}\le\varepsilon$, then
  $(\frac43-\varepsilon)\,\delta^{\mathrm{OT}}_{\mathrm{ed}}(G,H)\le\delta_{\mathrm{ed}}(G,H)$.
\item \lean{proposition6_delta} and \lean{theorem4_item1} (Theorem~3~(1)): for every $\varepsilon>0$ there is an odd
  $N\ge3$ with $(\frac43-\varepsilon)\,\delta^{\mathrm{OT}}_{\mathrm{ed}}(G,H)\le\delta_{\mathrm{ed}}(G,H)$.
\end{itemize}
Here $C_{\mathrm{ed}}$, $d_{\mathrm{ed}}$, $d^{\mathrm{OT}}_{\mathrm{ed}}$, $\delta_{\mathrm{ed}}=d_{\mathrm{ed}}/n^2$
and $\delta^{\mathrm{OT}}_{\mathrm{ed}}$ are \lean{costEditBij}/\lean{costEdit}, \lean{dEdit}, \lean{dOTEdit},
\lean{deltaEdit}, \lean{deltaOTEdit} of \file{Common/EditDistance.lean}, the same definitions as for
Theorem~6 and Theorem~15. $d_{\mathrm{ed}}$ is an infimum over bijections and $d^{\mathrm{OT}}_{\mathrm{ed}}$ an
infimum over $\mathrm{DS}(V,W)$.

\section{Where the proof lives}
\begin{longtable}{>{\raggedright\arraybackslash}p{5.2cm} >{\raggedright\arraybackslash}p{10.6cm}}
\toprule \textbf{File} & \textbf{Content} \\ \midrule\endhead \bottomrule\endfoot
\file{Construction.lean} & $G'$, $H'$ (core, markers, semicircles $I_r$), their degrees, the marker order, and
  $G$, $H$ by Lemma~22 with $\Delta=N+1$, $\tau=4N-1$, $\delta=1$ (App.~C.1, l.~677--694). \\
\file{Lemma28_MarkerCost.lean} & the cyclic distance $\ell$, compliant alignments, the cost split, Lemma~23. \\
\file{CoreEdges.lean} & counting the core-core pairs: $\mathrm{mixCore}=2(N-1)-2\,\mathrm{pres}$; averages of
  permutation matrices. \\
\file{Lemma29_FractionalCost.lean} & $\pi_\pm$, the witness $Q=\frac12(P_{\pi_-}+P_{\pi_+})$,
  $C_{\mathrm{ed}}(Q)=3N$, Lemma~24. \\
\file{Lemma30_BijectionCost.lean} & restriction to $V(G')$, displacements, Claim~26, Lemma~25. \\
\file{Proposition6.lean} & Proposition~5, its final computation, Theorem~3~(1). \\
\file{Common/Lemma27_LockingMarkers*.lean}, \file{Common/Claim26_DegreeBound.lean} & Lemma~22
  (locking markers) and Observation~21 (shared with Proposition~8). \\
\end{longtable}

\section{Proof layout}
\begin{center}
\begin{tikzpicture}
  \node[main] (t3) at (-4.5,0) {Theorem 3 (1)\\\lean{theorem4_item1}};
  \node[main] (p5) at (4.5,0) {Proposition 5\\\lean{proposition6_dEdit}\\\lean{proposition6}};
  \node[res] (calc) at (-4.5,-2) {final computation\\\lean{proposition6_delta}\\\lean{proposition6_delta_of_le}};
  \node[res] (l24) at (-4.5,-4.4) {Lemma 24\\\lean{lemma29}\\\lean{costEdit_witness}};
  \node[res] (l25) at (4.5,-4.4) {Lemma 25\\\lean{lemma30}, \lean{lemma30_bij}};
  \node[res] (core) at (-6,-7) {core-core pairs\\\lean{mixCore_eq}\\\lean{coreCost_witness}};
  \node[res] (l23) at (0,-7) {Lemma 23\\\lean{lemma28}\\\lean{costEdit_extMarkers}};
  \node[res] (c26) at (6,-7) {Claim 26\\\lean{claim31}\\\lean{sum_disp_ge}};
  \node[shared] (l22b) at (-3,-9.4) {Lemma 22 (2), $\diamondsuit=\mathrm{ed}$\\\lean{lemma27_cost_edit}};
  \node[shared] (l22a) at (3.5,-9.4) {Lemma 22 (1), $\diamondsuit=\mathrm{ed}$\\\lean{lemma27_bij_edit}};
  \node[shared] (o21) at (3.5,-11.4) {Observation 21\\\lean{claim26_bij}};
  \draw[dep] (t3) -- (calc);
  \draw[dep] (calc) -- (l24);
  \draw[dep] (calc) -- (l25);
  \draw[dep] (p5) -- (l24);
  \draw[dep] (p5) -- (l25);
  \draw[dep] (l24) -- (core);
  \draw[dep] (l24) -- (l23);
  \draw[dep] (l25) -- (l23);
  \draw[dep] (l25) -- (c26);
  \draw[dep, rounded corners] (l25.south east) -- (8.2,-5.2) -- (8.2,-9.4) -- (l22a.east);
  \draw[dep, rounded corners] (l25.south west) -- (2.6,-5.6) -- (-3.2,-5.6) -- (core.north east);
  \draw[dep] (l23) -- (l22b);
  \draw[dep] (l22a) -- (o21);
\end{tikzpicture}
\end{center}
Arrows mean ``uses''. Blue: this folder; grey: \file{GsimOT/Common}. Nothing is cited; all graphs are
built in \file{Construction.lean}, which uses \lean{lockGraph} and \lean{lockSize} of Lemma~22.

\section{The proof step by step}

\paragraph{Construction (App.~C.1, pp.~21--22, l.~677--694; Figures 1, 4).} The graphs $G'$, $H'$ share the
markers $z_0,\dots,z_{N-1}$; the cores are $S_G=\{v_0,v_2,\dots,v_{2N-2}\}$ and $S_H=\{w_1,\dots,w_{2N-1}\}$,
placed at their indices on $\mathbb Z_{2N}$. Distinct core vertices are adjacent iff their positions sum to $0$
mod $2N$, and $z_r$ is adjacent to the core vertices on $I_r=\{(r+s)\bmod 2N\mid s<N\}$. Lean: \lean{Base N}
$=$ \lean{Fin N} $\oplus$ \lean{Fin N} is the common vertex set of $G'$ and $H'$: \lean{Sum.inl c} is $v_{2c}$ in
$G'$ and $w_{2c+1}$ in $H'$ (the paper's identification $v_i\leftrightarrow w_{i+1}$), \lean{Sum.inr r} is $z_r$;
\lean{posG}, \lean{posH}, \lean{semicircle}, \lean{coreGraph}, \lean{Gp}, \lean{Hp}. For odd $N$ the core edges
are $\{c,c'\}$ with $c+c'=N$ in $G'$ and with $c+c'+1=N$, $c\ne c'$ in $H'$ (\lean{Gp_adj_inl_inl},
\lean{Hp_adj_inl_inl}); $v_0$ and $w_N$ are isolated (l.~692--693, 722--723).

The graphs $G,H$ are obtained by Lemma~22 with the markers ordered by their degree in $G'$ (\lean{markerEnum},
via \lean{Tuple.sort}; \lean{markerEnum_mono}) and $\Delta=N+1$, $\tau=4N-1$, $\delta=1$ (l.~688--690):
\lean{GN N} $=$ \lean{lockGraph (Gp N) (markerEnum N) (N + 1) (4 * N - 1) 1}, likewise \lean{HN N}, and the
clique has \lean{mK N} $=$ \lean{lockSize N (N+1) (4N-1) 1} $=2\Delta+2\tau+k(\tau+\delta)+1$ vertices.
The paper's validity check of the parameters (l.~690--691) is \lean{deg_Gp_le}, \lean{deg_Hp_le}
(all degrees $\le N+1$) and \lean{deg_Gp_inr}, \lean{deg_Hp_inr}, \lean{markerEnum_deg} (the degrees of $z_r$
differ by one); these are exactly the hypotheses of \lean{lemma27_bij_edit}.

\paragraph{Compliant alignments and $\ell$ (l.~695--699).} A compliant $Q$ is \lean{extendId (extMarkers R)}:
$R$ on the core vertices, extended by the identity on the markers (\lean{extMarkers}) and on the clique
(\lean{extendId}, Common). $\ell(v_i,w_j)=\min_k|i-j+2kN|$ is \lean{cycDist}/\lean{ell}.

\paragraph{Lemma 23 (p.~22, l.~700--706).} Mapping $v_i$ to $w_j$ changes the adjacency to exactly
$\ell(v_i,w_j)$ markers: \lean{card_semicircle_ne} (the number of semicircles containing exactly one of two
positions is their cyclic distance), summed in \lean{sum_marker_inl}, \lean{sum_marker_inr}. Counting both
orientations, \lean{costEdit_extMarkers} splits $C_{\mathrm{ed}}$ w.r.t.\ $(G',H')$ into the core-core part
\lean{coreCost N R}, the core-marker part $2\sum R(c,d)\,\ell(c,d)$, and $0$ for marker pairs;
\lean{lemma28} transfers this to $(G,H)$ by Lemma~22~(2) (\lean{lemma27_cost_edit}).

\paragraph{Lemma 24 (pp.~22--23, l.~707--725).} $\pi_\pm(v_i)=w_{i\pm1}$; with the identification, $\pi_+$ is the
identity on the core and $\pi_-$ is \lean{predCore} ($c\mapsto c-1\bmod N$). $Q=\frac12(P_{\pi_-}+P_{\pi_+})$ is
\lean{Qwit N} $=$ \lean{extendId (extMarkers (Rwit N))}, \lean{Rwit} $=$ \lean{avgPerm}; \lean{Qwit_isDS}.
Core-marker pairs: $\ell(v_i,\pi_\pm(v_i))=1$ (\lean{ell_self}, \lean{ell_predCore}), so they give $2N$
(\lean{sum_ell_witness}, l.~713--714). Core-core pairs (l.~715--724, Figure~5): \lean{coreCost_avgPerm} writes
them as the average of the four cases $(\pi_\pm,\pi_\pm)$, and \lean{mixCore_eq} gives each case as
$2(N-1)-2\,\mathrm{pres}$, where \lean{pres} counts preserved ordered core edges; \lean{pres_id_id},
\lean{pres_pred_pred} (no core edge preserved) and \lean{pres_id_pred}, \lean{pres_pred_id} ($N-2$ of the
$N-1$ ordered core edges preserved, the exceptional one of Figure~5b is lost) give $\frac14(2(N-1)+2+2+2(N-1))=N$ (\lean{coreCost_witness}). Hence
$C_{\mathrm{ed}}(Q)=3N$ (\lean{costEdit_witness}) and $d^{\mathrm{OT}}_{\mathrm{ed}}\le3N\le3N+1$ (\lean{lemma29},
via \lean{dOTEdit_le}).

\paragraph{Lemma 25 (pp.~23--24, l.~726--753).} Let $\pi$ be a bijection. If $C_{\mathrm{ed}}(\pi)\ge\tau=4N-1$, done
(the case split in \lean{lemma30_bij}). Otherwise Lemma~22~(1) (\lean{lemma27_bij_edit}) gives that $\pi$
fixes the markers and maps $V(G')$ onto itself (l.~727--728). The paper then replaces $\pi$ on the clique by
the identity (l.~728--729); Lean instead restricts $\pi$ to $V(G')$ (\lean{restrictInl}, applied twice to get
the core map $\rho$) and drops the pairs that involve the clique, which only lowers the cost
(\lean{costEditBij_restrict}). The remaining cost is \lean{mixCore N rho rho} plus
$2\sum_c\ell(v_{2c},\pi(v_{2c}))$ (\lean{costEditBij_sumCongr}, by Lemma~23). Displacements
$d_i\in\{-N+2,\dots,N\}$ (l.~729--732) are \lean{disp}; $\ell(v_i,\pi(v_i))=|d_i|$ (\lean{ell_eq_abs_disp}),
$|d_i|\ge1$ (\lean{one_le_ell}); $U=\sum|d_i|$ and $2k=$ \lean{pres N rho rho}.

\emph{Claim 26 (p.~24, l.~737--749)}, \lean{claim31}: $U\ge N+2k$. The case $d_i=N$ (l.~738--739) and the
general case follow the paper: $\sum d_i\equiv N\pmod{2N}$ (\lean{sum_disp}, l.~740--741), hence $|\sum d_i|\ge N$
(\lean{sum_disp_ge}, l.~742); a preserved edge has $d_i+d_{i'}=0$ (\lean{Preserved}, \lean{disp_add_disp_partner},
l.~743--745); the set $S$ of endpoints of preserved edges has $|S|=2k$ (\lean{pres_eq_card}) and
$\sum_{S}d_i=0$ is proved with the involution $c\mapsto N-c$ on $S$ (l.~746--749).

Finally (l.~750--753) the core-core pairs cost $2(N-1-2k)$ (\lean{mixCore_eq}, \lean{sum_adj_Gp_core},
\lean{sum_adj_Hp_core}) and $C_{\mathrm{ed}}(\pi)\ge2(N-1-2k)+2U\ge4N-2$ (\lean{lemma30_bij});
\lean{lemma30} takes the infimum (\lean{le_dEdit}).

\paragraph{Proof of Proposition 5 (p.~24, l.~754--758).} \lean{proposition6_delta_of_le} is the paper's chain
$d_{\mathrm{ed}}\ge4N-2=(\frac43-\frac{10}{3(3N+1)})(3N+1)\ge(\frac43-\frac{10}{3(3N+1)})\,d^{\mathrm{OT}}_{\mathrm{ed}}
\ge(\frac43-\varepsilon)\,d^{\mathrm{OT}}_{\mathrm{ed}}$ (a \texttt{calc} block, using \lean{lemma29}, \lean{lemma30}
and $\frac43-\frac{10}{3(3N+1)}\ge0$), followed by division by the common normalisation $n^2$ (l.~757).
\lean{proposition6_delta} picks the odd $N=2k+3$ with $k>\frac{10}{9\varepsilon}$, so that
$\frac{10}{3(3N+1)}\le\varepsilon$ (l.~757--758).

\section{Assumptions and trusted base}
No cited result and no hypothesis beyond $N$ odd, $N\ge3$ (and $\varepsilon>0$). The main theorems depend only
on \lean{propext}, \lean{Classical.choice}, \lean{Quot.sound} (checked by \file{scripts/CheckAxioms.lean}).

\section{Notes on the formalisation}
\begin{itemize}
\item \textbf{Statement.} \lean{proposition6_dEdit} is the paper's statement; \lean{proposition6} is
  the form with an explicit witness $Q$. The ``sufficiently large $N$'' of the second sentence is made explicit by
  \lean{proposition6_delta_of_le} (every odd $N\ge3$ with $\frac{10}{3(3N+1)}\le\varepsilon$).
\item \textbf{Lemma 24.} The Lean computes the cost of the witness, $C_{\mathrm{ed}}(Q)=3N$
  (\lean{costEdit_witness}), and \lean{lemma29} derives the bound $3N+1$ from it.
\item \textbf{Lemma 25.} Dropping the clique pairs instead of replacing $\pi$ on the clique by the identity
  (same effect: the cost only decreases). Claim~26 is stated with $2k=$ \lean{pres N rho rho}, the number of
  preserved \emph{ordered} core edges.
\item \textbf{Vertex sets.} $G$ and $H$ live on the same type \lean{VN N}; the identification
  $v_i\leftrightarrow w_{i+1}$ is built into \lean{Base N}.
\item \textbf{Lemma 22 for $\diamondsuit=\mathrm{ed}$} (\file{Common/Lemma27_LockingMarkersEdit.lean}): item~(1)
  is derived from the case $\diamondsuit=1$ via $C_1(\pi)\le C_{\mathrm{ed}}(\pi)$
  (\lean{costOneBij_le_costEditBij}); the paper proves both cases at once from Eq.~(16) and Observation~21.
\end{itemize}

\section{How to check this result}
\begin{itemize}
\item Build: \texttt{lake build GsimOT.Proposition6\_EditLowerBound.Proposition6} (or \texttt{bash build.sh}
  for the whole library, which also runs \file{scripts/CheckAxioms.lean}).
\item Axioms: \texttt{\#print axioms GsimOT.Prop6.proposition6\_dEdit}, and likewise for
  \lean{proposition6}, \lean{proposition6_delta}, \lean{proposition6_delta_of_le}, \lean{theorem4_item1};
  expected output: \lean{propext}, \lean{Classical.choice}, \lean{Quot.sound}.
\item The table \file{Proposition6_correspondence.pdf} in this folder lists every paper item with its
  Lean declaration and source location.
\end{itemize}
\end{document}
